% RACER-GT 預先註冊 protocol 範本（繁體中文）。
\documentclass[12pt,a4paper,oneside]{article}
% Traditional Chinese setup. Compile with XeLaTeX.
\usepackage{fontspec}
\usepackage{xeCJK}
\setmainfont{Noto Serif}
\setsansfont{Noto Sans}
\setmonofont{Noto Sans Mono}
\setCJKmainfont{Noto Serif CJK TC}
\setCJKsansfont{Noto Sans CJK TC}
\setCJKmonofont{Noto Sans CJK TC}
\XeTeXlinebreaklocale "zh"
\XeTeXlinebreakskip = 0pt plus 1pt

% Single source of truth for version and date across every RACER-GT document.
% Regenerate nothing by hand: bump here, rebuild, and every PDF agrees.
\newcommand{\racerversion}{1.1.0}
\newcommand{\racerdate}{2026-07-27}
\newcommand{\racerauthor}{Yi-Hao Lai}
\newcommand{\raceraffil}{Department of Finance, Da-Yeh University}
\newcommand{\raceremail}{yhlai@mail.dyu.edu.tw}
\newcommand{\racerrepo}{https://github.com/yhlai0911/RACER-GT}

% Shared package set and mathematical macros for every RACER-GT document,
% in both languages. Language-specific font and theorem-name setup lives in
% _preamble-zh.tex and _preamble-en.tex.
%
% Font policy: the build depends only on the five Noto families below. They are
% installable on Linux CI (fonts-noto-core, fonts-noto-cjk) and on macOS via
% Homebrew casks. Do not introduce a font outside this set without adding it to
% .github/workflows/docs.yml as well, or the PDF build stops being reproducible.

\usepackage[margin=2.35cm,headheight=15pt]{geometry}
\usepackage{amsmath,amssymb,amsthm,mathtools,bm}
\usepackage{booktabs,longtable,tabularx,array,multirow}
\usepackage{graphicx,float,caption,subcaption}
\usepackage{enumitem}
\usepackage{xcolor}
\usepackage{microtype}
\usepackage{setspace}
\usepackage{fancyhdr}
\usepackage{hyperref}
\usepackage[nameinlink,noabbrev]{cleveref}
\usepackage{algorithm}
\usepackage{algpseudocode}
\usepackage{listings}
\usepackage{tikz}
\usetikzlibrary{arrows.meta,positioning,shapes.geometric,fit,calc}

\hypersetup{
  colorlinks=true,
  linkcolor=black,
  citecolor=black,
  urlcolor=blue,
  pdfauthor={\racerauthor}
}

\pagestyle{fancy}
\fancyhf{}
\rhead{v\racerversion}
\cfoot{\thepage}
\setstretch{1.28}
\setlength{\parskip}{0.35em}
\setlist[itemize]{leftmargin=2em,itemsep=0.25em,topsep=0.3em}
\setlist[enumerate]{leftmargin=2.2em,itemsep=0.25em,topsep=0.3em}

% ---------------------------------------------------------------- mathematics
\newcommand{\E}{\mathbb{E}}
\newcommand{\Var}{\operatorname{Var}}
\newcommand{\Cov}{\operatorname{Cov}}
\newcommand{\tr}{\operatorname{tr}}
\newcommand{\diag}{\operatorname{diag}}
\newcommand{\argmin}{\operatorname*{arg\,min}}
\newcommand{\argmax}{\operatorname*{arg\,max}}
\newcommand{\one}{\bm{1}}
\newcommand{\Rset}{\mathbb{R}}
\newcommand{\R}{\mathbb{R}}
\newcommand{\plim}{\operatorname*{plim}}
\newcommand{\RACER}{\textsc{Racer-GT}}
\newcommand{\indic}[1]{\mathbb{I}\!\left\{#1\right\}}
\newcommand{\Sig}{\bm{\Sigma}}
\newcommand{\wvec}{\bm{w}}
\newcommand{\evec}{\bm{e}}
\newcommand{\Zvec}{\bm{Z}}

% Design facets, used identically in both languages so that a reader can move
% between the Chinese and English documents without re-learning notation.
\newcommand{\facetT}{\mathcal{T}}   % historical date
\newcommand{\facetD}{\mathcal{D}}   % collection day
\newcommand{\facetS}{\mathcal{S}}   % collection stream
\newcommand{\facetR}{\mathcal{R}}   % technical replicate

\lstset{
  basicstyle=\ttfamily\footnotesize,
  breaklines=true,
  frame=single,
  rulecolor=\color{gray!40},
  backgroundcolor=\color{gray!5},
  columns=fullflexible,
  keepspaces=true,
  showstringspaces=false
}


\setlength{\parindent}{2em}

\newtheorem{assumption}{假設}[section]
\newtheorem{proposition}{命題}[section]
\newtheorem{corollary}{推論}[section]
\newtheorem{remark}{說明}[section]
\newtheorem{definition}{定義}[section]
\newtheorem{lemma}{引理}[section]

\renewcommand{\proofname}{證明}
\renewcommand{\abstractname}{摘要}
\renewcommand{\contentsname}{目錄}
\renewcommand{\refname}{參考文獻}
\renewcommand{\tablename}{表}
\renewcommand{\figurename}{圖}

\crefname{assumption}{假設}{假設}
\crefname{proposition}{命題}{命題}
\crefname{corollary}{推論}{推論}
\crefname{definition}{定義}{定義}
\crefname{lemma}{引理}{引理}
\crefname{equation}{式}{式}
\crefname{table}{表}{表}
\crefname{figure}{圖}{圖}
\crefname{section}{節}{節}

\lhead{RACER-GT：Protocol 與預先註冊}

\title{\bfseries RACER-GT Protocol 與預先註冊範本\\[0.4em]
\large 在看到任何資料之前必須鎖定的內容}
\author{\racerauthor\\\small 大葉大學 財務金融學系}
\date{\racerdate\ \textbullet\ 版本 \racerversion}

\begin{document}
\maketitle

\begin{abstract}
\noindent 本範本應在蒐集開始\emph{之前}填寫、加上時戳並存放。其目的狹窄而明確：讓「取得端的研究者自由度」可被看見。一條在分析者已看過它將用來解釋的財金結果之後才建構出來的 Google Trends 數列，不是量測，而是選擇。以下所有內容，都是為了讓這個區別能被他人查核。
\end{abstract}

\tableofcontents
\newpage

\section{研究識別}

\begin{itemize}
\item 標題、作者、單位、聯絡方式。
\item 存放位置與時戳（OSF、AEA RCT Registry、機構典藏，或簽署的 Zenodo deposit）。
\item 使用的 RACER-GT 版本。
\item \textbf{Protocol hash}（由 \texttt{racergt design} 產生），請記錄於此。日後任何對鎖定設定的改動都會改變此 hash——這正是讓設定漂移「可被偵測」而非「可被否認」的機制。
\end{itemize}

\section{構念與查詢設定}

請先陳述經濟構念、再陳述查詢，順序不可顛倒；顛倒順序等於邀請自己挑選能得到理想結果的查詢。

\begin{itemize}
\item 本數列意圖量測的構念。
\item 為何此關鍵字或 Topic 能量測它，以及它還可能一併捕捉到什麼。
\item \texttt{keyword}、\texttt{topic\_or\_term}、\texttt{geo}、\texttt{category}、\texttt{search\_property}、\texttt{language}。
\item 類別選擇的理由。限制在 Finance 會改變構念，不限制同樣會改變。請說明您選了哪一個、為什麼。
\item 歷史起訖日期，以及正規化所用的 baseline 期間。
\end{itemize}

\section{蒐集設計}

\begin{itemize}
\item Collection-day offsets、streams、technical replicates、time slots。
\item 每個 stream 固定了什麼：裝置、瀏覽器 profile、cookie 狀態、帳號狀態、網路路徑、IP。
\item \textbf{聲明：}除非執行了「僅變動 IP」的 crossover factorial design，否則估出的 stream effect 不會被報告為 IP effect。
\item Chunk 視窗長度、step、最小重疊。
\item 蒐集途徑（官方 API、手動匯出、機構核准的 client）及其合規依據。
\end{itemize}

\section{事前鎖定的分析設定}

設定檔的每一個欄位都屬於 protocol。至少須陳述：校準門檻 $c_{\min}$、Huber 常數、參考 chunk 策略、正規化模式、duplicate 門檻、共變異數估計式、權重上限、G-study 轉換、benchmark 模式與權重。

\section{驗收門檻}

請在蒐集\emph{之前}填妥每一個門檻，並承諾無論結果為何都據實回報。

\begin{table}[H]
\centering\small
\begin{tabularx}{\textwidth}{lXr}
\toprule
規則 & 防範的問題 & 您的設定值\\
\midrule
\texttt{min\_unique\_pulls} & 批次大多是重複值 & \\
\texttt{min\_spectral\_effective\_pulls} & 把相依當成重複量測 & \\
\texttt{max\_zero\_share} & 稀有詞彙無法支撐日頻 & \\
\texttt{min\_level\_generalizability} & 水準無法重現 & \\
\texttt{min\_innovation\_generalizability} & 日變動無法重現 & \\
\texttt{min\_detection\_kappa} & 零／非零型態不穩定 & \\
\texttt{max\_component\_share} & 單一相依叢集主導全批 & \\
\texttt{max\_convergence\_mae\_100} & consensus 尚未收斂 & \\
\texttt{max\_benchmark\_standardized\_rmse} & 日頻與週頻數列不一致 & \\
\bottomrule
\end{tabularx}
\end{table}

\section{Pilot 分離}

明確說明是否使用初期監控窗，並承諾以下規則：若監控導致 threshold、stitching rule、query、code 或 decision tree 發生\emph{任何}改變，被監控的資料即成為 pilot 資料，正式批次須由新的 Day 0 重新開始。缺少這項承諾，同一批資料就同時被用來設計與驗證驗收規則，產生的 acceptance bias 無法靠事後謹慎補救。

\section{下游使用}

\begin{itemize}
\item 此指標將進入的財金模型，現在就要指定。
\item 以水準或差分形式進入。若為差分，具約束力的是 innovation generalizability 門檻，而非 level 門檻。
\item 測量誤差的處理方式：reliability-corrected OLS、SIMEX，或兩者並用。
\item \textbf{聲明：}GT 處理規則在檢視任何報酬、波動、價差或 VaR 結果之前即已鎖定。
\end{itemize}

\section{偏離紀錄}

保留一節於事後填寫，記錄對本 protocol 的每一項偏離、決定的時點與理由。一份有誠實偏離紀錄的 protocol，比一份完全沒有偏離的更可信。

\end{document}
